% secciones/introduccion.tex -- Sec 1 Introducción

La bioconversión de residuos orgánicos con \textit{Hermetia illucens} se ha consolidado como una vía de valorización: transforma desperdicios agroindustriales en biomasa rica en proteína para alimentación animal y un residuo estabilizado que devuelve materia orgánica al suelo. Su evaluación ambiental completa, sin embargo, exige contabilizar todo el carbono del sistema, incluido el que retorna a la atmósfera durante el proceso. En un contenedor de cría, la emisión de CO$_2$ medida mezcla dos respiraciones que ningún sensor separa: la de las larvas y la de la microbiota que degrada el alimento no consumido. La distinción no es un detalle para el balance de gases de efecto invernadero del proceso: la respiración larvaria es un costo metabólico del insecto productivo, mientras que la del lecho depende del manejo del sustrato; y atribuir el total medido a cada fuente no puede lograrse con la medición sola, sino que requiere un modelo mecanístico de ambas.

El modelo dinámico de \textcite{eriksenDynamicModellingFeed2022} es hoy la descripción más completa de la respiración larvaria: deriva la producción de CO$_2$ de la fisiología del insecto ---asimilación, crecimiento y acumulación de lípidos--- con cierre de carbono, y reproduce trayectorias de emisión bajo las condiciones en que fue desarrollado. Sin embargo, esas reproducciones emplean los mismos escenarios con que se parametrizó, por lo que no constituyen una validación independiente; y el modelo no representa ni la degradación microbiana del sustrato ni la limitación por disponibilidad de alimento, factores ineludibles en un lote real con dietas variables. Los estudios de emisión en reactores con \textit{H. illucens} confirman que la señal de cámara mezcla ambas fuentes \parencite{ermolaevGreenhouseGasEmissions2019,parodiBioconversionEfficienciesGreenhouse2020a}.

El objetivo de este estudio es formular y evaluar un modelo de caja gris que conserve la estructura del núcleo fisiológico de \textcite{eriksenDynamicModellingFeed2022} y añada explícitamente la relación alimento--larva--lecho, de modo que las emisiones medidas en cámaras estáticas puedan atribuirse a sus dos fuentes. Tres preguntas guían la evaluación: (i) cuánto del error de predicción proviene de aplicar el núcleo original fuera de su dominio de calibración; (ii) qué aporte representa la fuente microbiana del lecho en el total de cámara; y (iii) si la calidad y la disponibilidad del alimento explican las diferencias entre dietas sin sustituir los balances respiratorios del modelo de referencia. La evaluación distingue transferencia sin reajuste, recalibración con datos propios y comparación entre configuraciones anidadas, e incluye validación cruzada interna y un análisis de incertidumbre con límites declarados. La contribución del estudio es la formulación de la fuente microbiana del lecho y del acople entre la disponibilidad de alimento y la asimilación larvaria ---ambos ausentes en el modelo de referencia---, junto con un protocolo de comparación que distingue la recalibración del núcleo de la incorporación conjunta de las extensiones, sin atribuir de forma independiente la mejora a mecanismos que las configuraciones no aíslan.
